\chapter{Java}
\label{cap:java}

\begin{edital}
Java (versão 6+), JavaEE (6+), JakartaEE, JPA (2+), frameworks JUnit,
Hibernate, JSF, Primefaces, Spring/SpringCloud/SpringBoot. É a linguagem-base
do perfil.
\end{edital}
\peso{MÉDIO --- na amostra de provas, Java concentrou-se no TJ-RJ (Analista de
Sistemas); caiu pouco na Dataprev 2024. Mas está no edital: estude os pares
que a FGV adora inverter.}

\section{Exceções: checked $\times$ unchecked}

\begin{center}
\begin{tabular}{@{}p{2.6cm}p{4.6cm}p{4.6cm}@{}}
\toprule
& \textbf{Checked (verificada)} & \textbf{Unchecked (não verificada)} \\
\midrule
Superclasse & \texttt{Exception} (não \texttt{RuntimeException}) & \texttt{RuntimeException} \\
Compilador & \textbf{exige} try-catch \textbf{ou} \texttt{throws} & não exige \\
Exemplos & \texttt{IOException}, \texttt{SQLException} & \texttt{NullPointerException}, \texttt{ArrayIndexOutOfBounds}, \texttt{IllegalArgument} \\
\bottomrule
\end{tabular}
\end{center}

\begin{pegadinha}
Pegadinha nº 1 do bloco: trocar os exemplos. \texttt{NullPointerException} é
\textbf{unchecked} (estende \texttt{RuntimeException}); \texttt{IOException}
é \textbf{checked}. Absoluto a descartar: ``unchecked exige throws
obrigatório'' --- é o contrário.
\end{pegadinha}

\section{Polimorfismo: overload $\times$ override}

\begin{center}
\begin{tabular}{@{}p{2.6cm}p{4.6cm}p{4.6cm}@{}}
\toprule
& \textbf{Overload (sobrecarga)} & \textbf{Override (sobrescrita)} \\
\midrule
Onde & mesma classe & subclasse \\
Assinatura & \textbf{muda} (parâmetros diferentes) & \textbf{igual} à do pai \\
Resolução & compile-time (estática) & runtime (dinâmica) \\
Anotação & --- & \texttt{@Override} \\
\bottomrule
\end{tabular}
\end{center}

\begin{pegadinha}
Overload = mesmo nome, assinatura diferente; override = mesma assinatura na
subclasse. Sobrescrever \textbf{não} viola Liskov por si só.
\end{pegadinha}

\section{OO e SOLID}

Pilares: abstração, encapsulamento, herança, polimorfismo.

\begin{center}
\begin{tabular}{@{}p{1.2cm}p{9.6cm}@{}}
\toprule
\textbf{Letra} & \textbf{Princípio} \\
\midrule
S & Single Responsibility --- uma responsabilidade por classe \\
O & Open/Closed --- aberta a extensão, fechada a modificação \\
L & \textbf{Liskov} --- subtipo substitui o tipo base sem quebrar o programa \\
I & Interface Segregation --- interfaces pequenas e específicas \\
D & Dependency Inversion --- depender de abstração, não de implementação \\
\bottomrule
\end{tabular}
\end{center}

\begin{pegadinha}
Não confundir \textbf{ISP} (interface enxuta) com \textbf{SRP} (uma
responsabilidade) --- são os dois princípios que a FGV mais troca no SOLID.
Liskov: a subclasse pode sobrescrever, desde que continue substituível.
\end{pegadinha}

\section{Coleções (Collections)}

\begin{center}
\begin{tabular}{@{}p{2.2cm}p{3.8cm}p{4cm}@{}}
\toprule
\textbf{Interface} & \textbf{Implementações} & \textbf{Característica} \\
\midrule
\texttt{List} & \texttt{ArrayList}, \texttt{LinkedList} & ordenada, permite duplicatas \\
\texttt{Set} & \texttt{HashSet}, \texttt{TreeSet}, \texttt{LinkedHashSet} & sem duplicatas \\
\texttt{Map} & \texttt{HashMap}, \texttt{TreeMap}, \texttt{LinkedHashMap} & chave-valor \\
\bottomrule
\end{tabular}
\end{center}

\textbf{ArrayList} (array dinâmico, acesso O(1) por índice) $\times$
\textbf{LinkedList} (lista ligada, inserção/remoção O(1) no meio se já tem o
nó). \texttt{HashMap} (sem ordem) $\times$ \texttt{LinkedHashMap} (ordem de
inserção) $\times$ \texttt{TreeMap} (ordenado). \texttt{getOrDefault} retorna
o \textbf{default} se a chave não existe --- não lança exceção nem retorna
erro HTTP. \texttt{==} compara \textbf{referência}; \texttt{.equals()}
compara \textbf{conteúdo}.

\section{JVM, JPA e JavaEE}

\begin{itemize}
\item \textbf{JVM:} executa bytecode; \textbf{heap} (objetos),
\textbf{garbage collection} (libera memória sem referência); ferramentas de
monitoramento (\texttt{jconsole}, \texttt{jps}, \texttt{jstack}).
\item \textbf{JPA} (especificação de persistência) $\times$
\textbf{Hibernate} (implementação ORM). \texttt{EAGER} $\times$
\texttt{LAZY} (carregamento antecipado $\times$ sob demanda). Problema
\textbf{N+1} (uma consulta por relação): resolve com \texttt{LAZY} +
\texttt{JOIN FETCH} na consulta.
\item \textbf{JavaEE/JakartaEE:} EJB, JPA, JMS, JSF; \textbf{Jakarta} é a
continuação do JavaEE sob a Eclipse Foundation (namespace \texttt{jakarta.*}
em vez de \texttt{javax.*}).
\item \textbf{JSF/Primefaces:} framework de UI server-side baseado em
componentes.
\end{itemize}

\begin{pegadinha}
JPA N+1: o distrator oferece EAGER em tudo ou ``deixar o JPA decidir'' como
solução --- a resposta certa é \textbf{LAZY + JOIN FETCH} na consulta que
precisa da relação.
\end{pegadinha}

\section{Recursos modernos (Java 17/21 LTS)}

\begin{itemize}
\item \textbf{Sealed / non-sealed / final:} uma classe \textbf{sealed}
exige subclasses declaradas (\texttt{permits}, ou no mesmo arquivo), e cada
subtipo deve ser \texttt{final}, \texttt{sealed} ou \texttt{non-sealed}.
\texttt{non-sealed} reabre a hierarquia.
\item \textbf{Threads virtuais (Project Loom)} $\times$ threads de
plataforma: milhares de \textbf{virtuais} compartilham uma \textbf{carrier
thread} de plataforma (não é uma-para-uma). Brilham em cenários I/O-bound.
\item \texttt{record}, \texttt{var}, switch expressions, pattern matching.
\item \textbf{Streams e lambdas} (Java 8+): programação funcional em
coleções.
\end{itemize}

\begin{pegadinha}
Sealed: o distrator diz que \texttt{final} ``não pode estender sealed''
(pode) ou inventa erro de escopo; o erro real costuma ser uma sealed
\textbf{sem nenhuma subclasse permitida}. Threads virtuais: o distrator diz
que cada virtual equivale a uma de plataforma (consumo linear) ou que o
número de threads de plataforma limita o de virtuais --- é o contrário:
muitas virtuais, poucas carriers.
\end{pegadinha}

\section{Leitura ativa de código em Java}
\label{sec:leitura-ativa-java}

Java é a linguagem onde a FGV mais gosta de esconder o erro \textbf{dentro}
de uma estrutura que parece correta à primeira vista --- um \texttt{try-catch}
que compila mas nunca captura a exceção certa, um escopo de variável que
parece visível mas não é, um \texttt{synchronized} que quebra a promessa da
thread virtual. Aplique o protocolo de leitura ativa do Capítulo
\ref{cap:programacao} (Seção \ref{sec:leitura-ativa-programacao}): releia
linha a linha, sem assumir.

\subsection{Checked/unchecked disfarçadas em try-catch}

\begin{conceito}[O que a FGV esconde no bloco try-catch]
\begin{itemize}
\item \textbf{Ordem dos catch:} um \texttt{catch (Exception e)} \textbf{antes}
de um \texttt{catch (IOException e)} \textbf{não compila} --- o catch mais
genérico precisa vir \textbf{por último}. A FGV mostra um trecho ``quase
certo'' com a ordem trocada e pergunta o que acontece: é \textbf{erro de
compilação}, não exceção não tratada em runtime.
\item \textbf{Catch que não cobre o tipo lançado:} um método declara
\texttt{throws SQLException} mas o bloco só tem
\texttt{catch (RuntimeException e)} --- isso \textbf{não compila}, porque
\texttt{SQLException} é \textbf{checked} e exige tratamento explícito ou
propagação com \texttt{throws} no método chamador.
\item \textbf{finally que engole a exceção:} um \texttt{return} dentro do
bloco \texttt{finally} \textbf{sobrescreve} qualquer exceção lançada no
\texttt{try}/\texttt{catch} --- a exceção original se perde silenciosamente.
É pegadinha clássica de ``qual valor o método retorna''.
\item \textbf{Multi-catch:} \texttt{catch (IOException | SQLException e)}
exige que as duas exceções \textbf{não tenham relação de herança entre si}
(uma não pode ser subtipo da outra) --- senão não compila.
\end{itemize}
\end{conceito}

\begin{pegadinha}
O distrator típico apresenta um bloco try-catch com a ordem dos \texttt{catch}
invertida (genérico antes do específico) e afirma que ele ``compila e
funciona normalmente'' --- \textbf{falso}, é erro de compilação. Outro
distrator: dizer que uma exceção \textbf{checked} lançada dentro de um método
\texttt{void} sem \texttt{throws} e sem try-catch ``é ignorada em runtime''
--- também falso, isso \textbf{não compila}. Erro de compilação e exceção em
runtime são coisas diferentes; a FGV testa se você distingue as duas.
\end{pegadinha}

\subsection{Escopo de variável --- a pegadinha silenciosa}

\begin{itemize}
\item \textbf{Shadowing (sombreamento):} uma variável local com o
\textbf{mesmo nome} de um atributo da classe \textbf{esconde} o atributo
dentro do método --- \texttt{this.atributo} continua acessível, mas
\texttt{atributo} sozinho refere-se à variável local. Enunciado que usa o
mesmo nome dos dois lados é sinal de alerta.
\item \textbf{Variável efetivamente final em lambda/classe anônima:} uma
variável local usada dentro de uma lambda ou classe anônima precisa ser
\textbf{final ou efetivamente final} (nunca reatribuída depois da
inicialização) --- código que reatribui essa variável depois de criar a
lambda \textbf{não compila}.
\item \textbf{Escopo de bloco em \texttt{for}:} uma variável declarada
\textbf{dentro} do \texttt{for} (\texttt{for (int i = 0; ...)}) não existe
fora dele --- usá-la depois do laço é erro de compilação, não um valor
residual do último loop.
\end{itemize}

\begin{pegadinha}
A FGV declara uma variável dentro de um bloco (\texttt{if}, \texttt{for},
\texttt{try}) e depois faz o código ``usá-la'' fora do bloco na alternativa
--- isso é erro de compilação (variável fora de escopo), não um valor
qualquer. Não confunda com JavaScript, onde \texttt{var} vaza escopo de
função --- em Java o escopo é sempre de \textbf{bloco}.
\end{pegadinha}

\subsection{Armadilhas de threads virtuais}

\begin{conceito}[Pinning --- quando a thread virtual não é tão leve assim]
Uma thread virtual roda sobre uma \textbf{carrier thread} de plataforma. Se o
código dentro de um bloco \texttt{synchronized} faz uma chamada
\textbf{bloqueante} (I/O, banco de dados), a thread virtual pode ficar
\textbf{presa (pinned)} à carrier em vez de liberá-la --- isso derruba
justamente a vantagem de escala das threads virtuais, porque agora uma
operação bloqueante ocupa uma carrier inteira. Na prática, o cenário de risco
é: \texttt{synchronized} + chamada bloqueante \textbf{dentro} do bloco
sincronizado, sob alta concorrência.
\end{conceito}

\begin{pegadinha}
O distrator diz que threads virtuais ``eliminam totalmente qualquer
preocupação com bloqueio ou concorrência'' --- \textbf{falso}: o cenário
clássico de \texttt{synchronized} + I/O bloqueante ainda pode prender a
thread virtual à sua carrier, anulando o ganho de escala. Outro distrator
troca a causa: dizer que o problema é ``excesso de memória'' (é
concorrência/pinning, não memória). Threads virtuais reduzem o \textbf{custo
de criar} threads, não eliminam todo cuidado com seções críticas.
\end{pegadinha}

\begin{jacaiu}
Checked $\times$ unchecked; overload $\times$ override; Liskov (SOLID);
coleções (\texttt{ArrayList}/\texttt{LinkedList}/\texttt{HashMap});
\texttt{==} $\times$ \texttt{equals}; sealed classes; threads virtuais e
pinning; JPA N+1 (TJ-RJ); leitura de REST controller
(\texttt{@RestController}, \texttt{@GetMapping}, \texttt{@PathVariable}) e
comportamento de \texttt{Map}/\texttt{HashMap} (\texttt{getOrDefault});
Spring Cloud Eureka (atributo \texttt{lease-expiration}); Hibernate Envers
(auditoria via Revision Listener); ordem de catch e exceção engolida por
\texttt{finally}; escopo de variável em bloco. Rode \texttt{./quiz.py java}.
\end{jacaiu}

\begin{comosair}
Par para decorar: checked (\texttt{Exception}, verificada em compilação,
trata ou declara \texttt{throws}) $\times$ unchecked (\texttt{RuntimeException},
não obriga tratamento); overload (mesmo nome, assinatura diferente,
compile-time) $\times$ override (mesma assinatura na subclasse, runtime,
\texttt{@Override}). N+1: LAZY para não carregar cedo + JOIN FETCH quando
precisar da relação numa consulta só. Em leitura de código, execute
mentalmente linha a linha e confie na semântica da API (\texttt{getOrDefault}
$\to$ default). Gatilhos de distrator: ``sempre'', ``independentemente da
carga'', ``equivale a uma de plataforma''. Nomes de atributo/config (Eureka
\texttt{lease-expiration}, Envers Revision Listener): a FGV troca por nomes
plausíveis (\texttt{heartbeat-interval}, interceptor, filtro) --- decore o
nome exato.
\end{comosair}
